\section{Appendix}\label{section_appendix}
We provide in this Appendix illustrations of the frequency and time focus functions introduced in the core of the paper. We stress that these do not intend to address specific applied problems, but simply to show  that such focus functions can indeed be designed and achieve well targeted goals.
\subsection{An example of frequency focus function}
\label{sec:app.freq.foc.example}
We revisit here Example~\ref{ex:frequency_focus}, and illustrate the behaviour of a frequency focus function similar to~\eqref{fo:ex.freq.foc} on a toy example of harmonic signal, involving sine waves with close frequencies in the low frequency range, and sine waves with better separated frequencies in the high frequency range. A spectrogram (square modulus of STFT) with suitably chosen analysis window is displayed in Figure~\ref{fi:freq.foc.spgram}. The spectrogram was computed using LTFATpy~\cite{ltfatpy}, the Python port of the Large Time-Frequency Analysis Toolbox~\cite{ltfat}. Clearly, the window choice for the STFT allows separation of sine waves in the high frequency part of the spectrogram, while in the low-frequency range these sine waves are mixed by the window, resulting in beating phenomenon (oscillations of the modulus due to interference between sine waves). Clearly enough, these two different behaviours result from a suitable choice of the STFT window scale.

\begin{figure}[h!]
    \centering
    \includegraphics[scale=.6]{freq_foc_spgram}
    \caption{Spectrogram of a harmonic signal involving sine waves with close frequencies in the low frequency range, and sine waves with better separated frequencies in the high frequency range.}
    \label{fi:freq.foc.spgram}
\end{figure}

\begin{figure}[h!]
    \centering
    \includegraphics[scale=.6]{freq_foc}
    \caption{Frequency focus function based upon a ratio $L^2/L^1$ norms, for the signal whose spectrogram is displayed in Figure~\ref{fi:freq.foc.spgram}}
    \label{fi:freq.foc}
\end{figure}
The prototypical frequency focus function introduced in~\eqref{fo:ex.freq.foc} turns out to be insufficient to produce the expected behaviour, because of spurious oscillations (not shown here). The latter can be suppressed by appropriate low-pass filtering. After such filtering and normalisation, on obtains a focus function as given in Figure~\ref{fi:freq.foc}. The latter can be suitably rescaled to match a prescribed range, and thresholded to achieve the support property assumed in its definition. Since it is not the objective of this paper to provide realistic choices for applications, we d'on't  address such a fine tuning here.

As can be seen on Figure~\ref{fi:freq.foc}, the frequency focus function takes lower values in the low-frequency part of the signal spectrum. This implies that the window in the frequency-focused transform will have larger scale, therefore smaller frequency bandwidth, which will provide the desired adaptation to the signal if the aim is to better separate sine waves.

\vspace{-5mm}
\begin{remark}
Other more explicit illustrations can also be obtained since all calculations rely on Gaussian integrals which are tractable. In our opinion this does not bring more insight than the numerical examples in the Appendix, we then refrain from following this avenue.
\end{remark}
\subsection{An example of time focus function}
\label{sec:app.time.foc.example}
We now revisit Example~\ref{ex:time_focus}, and illustrate the behaviour of a time focus function similar to~\eqref{eq_def2_sigma_tau} on another toy example, consisting in coloured noise with additional randomly located spikes with random amplitude. A realisation of such a random signal (involving 10 random peaks, out of which some are not visible because they are hidden by the background noise) is plotted in  Figure~\ref{fi:time.foc.sig}. Its spectrogram (not shown here) exhibits a clear frequency decay (as a result of the choice of noise) away from the spike locations (which do not decay in the frequency domain). The time focus function used here is designed to capture such a behaviour, which we illustrate below.

\begin{figure}
    \centering
    \includegraphics[scale=.6]{time_foc_sig.eps}
    \caption{Gaussian coloured noise with randomly located spikes, with random amplitudes.}
    \label{fi:time.foc.sig}
\end{figure}

\begin{figure}
    \centering
    \includegraphics[scale=.6]{time_foc.eps}
    \caption{Time focus function for the spike+noise signal plotted in Figure~\ref{fi:time.foc.sig}.}
    \label{fi:time.foc}
\end{figure}
We display in Figure~\ref{fi:time.foc} the corresponding time focus function, as defined in~\eqref{eq_def2_sigma_tau}. The focus function presents clear maxima near 7 time instants, which turn out to coincide with the locations of the 7 neatly visible spikes (out of 10 generated spikes, the other 3 having smaller amplitude). The function was rescaled to achieve the prescribed minimal value of 1; again further post-processing is possible to achieve the support condition of the definition, we refrained from doing it as this would require an extra thresholding parameter, such a fine tuning being irrelevant at this stage.

